% SPDX-FileCopyrightText: 2016-2026 Lari Natri <lari.natri@iki.fi>
% SPDX-License-Identifier: GPL-3.0-or-later
%
% Shared visual policy.
% This file is part of the 'ulsbs' package.
%
% Loaded immediately after ulsbs-typography.code.tex, before front matter.
% Requires imported graphics/color/header packages, defaults, book
% strings and typography. Owns common colors, font hooks, header styles,
% inline boxes and break policy; loads chapter marks at their original site.

\ProvidesFile{ulsbs-style.code.tex}[Shared visual policy]

% This is the final common policy layer, not the pre-typography bootstrap.
% For providecommand font/header defaults, consumers may define before loading
% or renew afterward. Hooks explicitly renewed here must be renewed
% after loading; changing one does not rerun the earlier bootstrap measurements.
% Use \definecolor/\colorlet after loading to replace named colors, \setlength
% for allocated lengths/skips, and direct assignment for penalty count registers.

% === Paragraph stretching and named colors ===

% Allow further emergency stretching to fix some "overfull hbox" situations.
% See: https://www.tex.ac.uk/FAQ-overfull.html
\setlength{\emergencystretch}{3em}

% color: Color definitions:
\definecolor{ulChordColor}{rgb}{0.20,0,0.22} % chords
\definecolor{ulTransposedKeyColor}{rgb}{0.60,0,0.04} % non-original prelude keys
\definecolor{ulAltInlineChordColor}{rgb}{0.4,0,0.22} % alt. chord coloring used with \ac
\definecolor{ulAltChordColorA}{rgb}{0.10,0.4,0.10} % green, set 1 (default)
\definecolor{ulAltChordColorB}{rgb}{0.35,0.22,0.06} % brown, set 2
\definecolor{ulAltChordColorC}{rgb}{0.17,0.23,0} % dirty green, set 3
\definecolor{ulAltChordColorD}{rgb}{0.13,0,0.42} % bluish, set 4
\definecolor{ulAltChordColorE}{rgb}{0.09,0.33,0.34} % cyan, set 5
\definecolor{ulAltChordColorF}{rgb}{0.36,0.15,0.23} % redish, set 6
\definecolor{ulMeasureBarColor}{rgb}{0.36,0,0.48} % measure bar
% encircled melody note color:
\definecolor{ulMelodyNoteColor}{rgb}{0.0,0.1536,0.1824}
% encircled melody note color (alt.):
\definecolor{ulMelodyNoteAltColor}{rgb}{0.2,0,0}
% playout mark color:
\definecolor{ulPlayoutMarkColor}{rgb}{0,0.5,0}
% color for repetition signs:
\definecolor{ulRepetitionSignColor}{rgb}{0,0,0.5}
\definecolor{ulSongNumberBackgroundColor}{rgb}{1,1,1} % song number background
\definecolor{ulNoteBackgroundColor}{rgb}{0.9,0.9,0.9} % textnote background
\definecolor{ulIndexBackgroundColor}{rgb}{1,1,1} % index alphabet background
\definecolor{ulInternalLinkColor}{rgb}{0,0,0.3} % internal links
\definecolor{ulUrlLinkColor}{rgb}{0.3,0,0.3} % external links
\definecolor{ulLinkBoxColor}{rgb}{0.5,0.9,0.6} % boxed links (view), not used now
% ulsbs@langsect-derived environments: the color for \up marks
\colorlet{ulLanguageSectionMarkColor}{ulChordColor}%
% ulsbs@langsect-derived environments: colors used for different languages
\definecolor{ulLanguageEnglishColor}{rgb}{0,0,0}%
\definecolor{ulLanguagePortugueseColor}{rgb}{0,0.26,0.26}%
\definecolor{ulLanguageFinnishColor}{rgb}{0,0.05,0.38}%

% === PDF presentation and preamble-finalized metadata ===

\hypersetup{linkbordercolor=ulLinkBoxColor,%
            linkcolor=ulInternalLinkColor,%
            urlcolor=ulUrlLinkColor,%
            unicode=true,%
            % TwoPageRight: display two pages side-by-side with flipping pages,
            %               odd pages on the right
            % TwoColumnRight: display two pages side-by-side with continuous
            %                 scrolling, odd pageson the right
            pdfpagelayout=TwoColumnRight,%
            pdfpagemode=UseNone,% alt. UseThumbs or UseOutlines (bookmarks)
            pdfprintscaling=None%
           }

\ExplSyntaxOn
% Defer book-string expansion until after author preamble configuration, before
% hyperref's begin-document metadata publication. Explicit hypersetup keys win,
% including empty values. Observe the existing handlers rather than replacing
% hyperref's parsing. Nonempty values from an earlier hyperref load also win.
\clist_map_inline:nn {title,author,subject,keywords}
  {
    \tl_if_empty:cF {@pdf#1} {\cs_gset:cpn {ulsbs@pdf@#1@explicit} {}}
    \exp_args:Nc \pretocmd {KV@Hyp@pdf#1}
      {\cs_gset:cpn {ulsbs@pdf@#1@explicit} {}}
      {} {\msg_fatal:nnn {ulsbs} {variant} {Cannot~observe~PDF~metadata~key}}
  }
% pdfinfo may contain arbitrary author fields; retain those and add only our
% version field when the author has not supplied DocumentVersion explicitly.
\cs_new_protected:Npn \ulsbs@finalizemetadata
  {
    \cs_if_exist:cF {ulsbs@pdf@title@explicit}
      {\hypersetup{pdftitle={\ulsbs@variantpdftitle}}}
    \cs_if_exist:cF {ulsbs@pdf@author@explicit}
      {\hypersetup{pdfauthor={\ulBookAuthor}}}
    \cs_if_exist:cF {ulsbs@pdf@subject@explicit}
      {\hypersetup{pdfsubject={\ulBookSubject}}}
    \cs_if_exist:cF {ulsbs@pdf@keywords@explicit}
      {\hypersetup{pdfkeywords={\ulBookKeywords}}}
    \cs_if_exist:cF {HyInfo@Value@DocumentVersion}
      {\hypersetup{pdfinfo={DocumentVersion={\ulBookVersion}}}}
  }
\AddToHook{begindocument/before}{\ulsbs@finalizemetadata}
\ExplSyntaxOff

\renewcommand{\ulSongNumberBackgroundColor}{ulSongNumberBackgroundColor} % song-number background in prelude
\renewcommand{\ulNoteBackgroundColor}{ulNoteBackgroundColor} % \musicnote and \textnote background
\renewcommand{\ulIndexBackgroundColor}{ulIndexBackgroundColor} % index alphabet section background


% === Font hooks and spacing policy ===

% Set lyric font style:
\providecommand{\ulLyricFont}{\sffamily\large}
% Verse layout: use the public lyric-font policy.
\renewcommand{\ulLyricFontStyle}{\ulLyricFont} % lyrics font style

% Set chord font style:
\providecommand{\ulChordFont}{\sffamily\bfseries\small\color{ulChordColor}}
% Chord rendering: use the public chord-font policy.
\renewcommand{\ulPrintChord}[1]{\ulChordFont#1}


% Chord character adjustment
%
% Adjust the appearance of some characters in chord definitions. Currently makes
% parenthesis smaller and raises them upwards. This is done to get low hanging
% characters of our default font to not collide with the lyrics below.
%
% To turn this off, set \ulAdjustChordChars{false}
%
\newcommand{\ulsbs@smraise}[1]{\ifulAdjustChordChars\raise0.234ex\hbox{\scalebox{0.81}#1}\else#1\fi}
\renewcommand{\ulChordLocals}{%
  \catcode`(\active
  \catcode`)\active
}
\newcommand{\ulsbs@smraiseopenparen}{\ulsbs@smraise{(}}
\newcommand{\ulsbs@smraisecloseparen}{\ulsbs@smraise{)}}
\begingroup
  \ulChordLocals
  \gdef({\ulsbs@smraiseopenparen}%
  \gdef){\ulsbs@smraisecloseparen}%
\endgroup

% Meter font style (add color):
\renewcommand{\ulMeterFont}{\color{ulChordColor}\scriptsize\sffamily\upshape}
\renewcommand{\ulNoteFont}{\small} % text/music note font (make smaller)
\renewcommand{\ulIndexTitleFont}{\sffamily} % index title font
\renewcommand{\ulIndexLyricFont}{\sffamily\slshape} % alternate (lyric) index title font
\renewcommand{\ulIndexTagFont}{\sffamily\small\slshape} % tag index font
% Text/music notes are prose even when inserted within an active verse.

% Chord accidentals: use a real sharp from the Noto Symbol font.
\renewcommand{\ulSharpSymbol}{\textsuperscript{{\ulSymbolFont♯}}}%
%
% Scalar verse spacing defaults are installed later, after font configuration.
% Share the annotation baseline even for empty beat-only anchors. Retained ink
% still determines height: beat marks do not acquire a full chord/pitch row.
% Authors may replace this hook to supply each anchor's clearance. The retained
% line pass can lift neighboring rows to that minimum without rerunning the hook.
\renewcommand{\ulChordLineParams}{%
  \baselineskip=\ulChordLineBaselineSkip%
  \lineskiplimit=\ulChordLineLineskipLimit%
  \lineskip=\ulChordLineLineskip%
}
% Song layout uses these flexible skips around prelude/body/postlude.
\ulBeforePostludeSkip=0pt plus 0pt minus 12pt % before the rights/postlude text
\ulAfterPreludeSkip=0pt plus 8pt minus 6pt % between prelude and retained body

%\addto\captionsfinnish{\renewcommand\chaptername{}} % sets chapter name prefix to nothing when babel: finnish
\addto\captionsenglish{\renewcommand\chaptername{}} % sets chapter name prefix to nothing when babel: english
\addto\captionsenglish{\renewcommand\appendixname{}} % sets appendix name prefix to nothing when babel: english

% === Page headers (fancyhdr) ===
%
% Note that some pages (the first pages of TOC and chapters etc) automatically
% use the style 'plain'. If needed, this style can also be modified with
% \fancypagestyle{plain}
%

% The style for book, chapter and section titles in page headers
\providecommand{\ulHeaderTitleStyle}{%
  \rmfamily\small\itshape%
}

% The style for page numbers in page headers
\providecommand{\ulPageNumberStyle}{%
  \rmfamily\small%
}

% This gets rid of implicit uppercasing, as it creates problems with images
% in the chapter titles. Upcasing is now done differently and controlled with
% \ulUppercaseBookTitleInHeader{true|false},
% \ulUppercaseSectionTitleInHeader{true|false} and
% \ulUppercaseShortChapterTitle{true|false}. Mark definitions stay outside styles.
\renewcommand{\chaptermark}[1]{\markboth{#1}{}}
\renewcommand{\sectionmark}[1]{\markright{#1}}

% Returns the book's title formatted for headers. Used in fancyhdr style below.
\newcommand{\ulsbs@headerbooktitle}{%
  {%
    \ulHeaderTitleStyle%
    \ifulUppercaseBookTitleInHeader%
      \protect\MakeUppercase{\expandafter\ulThisBookTitle}%
    \else%
      \ulThisBookTitle%
    \fi%
  }%
}

% Returns the sections's title formatted for headers. Used in fancyhdr style
% below.
\newcommand{\ulsbs@headersectiontitle}{%
  {%
    \ulHeaderTitleStyle%
    \ifulUppercaseSectionTitleInHeader%
      \protect\MakeUppercase{\expandafter\rightmark}
    \else%
      \rightmark%
    \fi%
  }
}

% Returns the chapter's title formatted for headers. Used in fancyhdr style
% below.
\newcommand{\ulsbs@headerchaptertitle}{%
  % Note that chapter title's optional upcasing is done elsewhere in the
  % definition of \ulMainChapter, because it cannot be done here due to possible
  % chapter images).
  {\ulHeaderTitleStyle{\leftmark}}%
}

\newcommand{\ulsbs@headerchapterandsectiontitles}{%
  \rmfamily\small%
  \ulsbs@headerchaptertitle%
  \ulsbs@ifnowidth{\rightmark}{}{%
    \hspace{1em}---\hspace{1em}\ulsbs@headersectiontitle%
  }%
}

% SPDX-FileCopyrightText: 2016-2026 Lari Natri <lari.natri@iki.fi>
% SPDX-License-Identifier: GPL-3.0-or-later
%
% Chapter corner marks.
% This file is part of the 'ulsbs' package.
%
% Requires TikZ and style defaults.
% Loaded with @ a letter, before header-style setup; no song-lifecycle dependency.
% Owns mark dimensions, dashed flag and cached PGF geometry. The surrounding
% style updates chapter color/dashed state; this file installs no chapter hooks.
% Header/footer callers draw overlays for the page being built, not while song
% source is captured. Default geometry is computed at load time.

\ProvidesFile{ulsbs-chapter-marks.code.tex}[ULSBS chapter corner marks]

% Setup the colored chapter corner mark for all pages within ulMainChapters.
%
% \ulChapterCornerMarkOffset should be at least 2mm to leave room for imprecise
% cutting of the page. \ulChapterCornerMarkOffset + \ulChapterCornerMarkWidth should
% be less than 4 mm, so that the whole mark is left outside of printer margins
% when printing on home printer to a paper that requires no cutting.
\newdimen\ulChapterCornerMarkOffset\ulChapterCornerMarkOffset=2.5mm%
\newdimen\ulChapterCornerMarkWidth\ulChapterCornerMarkWidth=.954915mm% 2.5 * 0.618034^2 :)
\newdimen\ulChapterCornerMarkLength\ulChapterCornerMarkLength=\dimexpr\ulTopMargin+\ulHeadHeight-\ulChapterCornerMarkOffset\relax%
% Gap count: 3 gaps -> 7 segments -> gaps are segments 2, 4 and 6.
\newcount\ulChapterCornerMarkDashCount\ulChapterCornerMarkDashCount=3%

% Set this to true to have dashed chapter corner marks. The surrounding style's
% chapter setup resets it to false; this module only reads the flag when drawing.
\ulsbs@newbool\ifulChapterCornerMarkDashed\ulsbs@chapcornermarkdashedtrue\ulsbs@chapcornermarkdashedfalse\ulChapterCornerMarkDashed\ulsbs@chapcornermarkdashedfalse

\newdimen\ulsbs@ccm@halfw
\newdimen\ulsbs@ccm@arm
\newdimen\ulsbs@ccm@reach
\newdimen\ulsbs@ccm@start
\newdimen\ulsbs@ccm@end
\newcount\ulsbs@ccm@segments
\newcount\ulsbs@ccm@segment

% Boundary #1 on either leg. Compute from the total so rounding never leaves
% a stray bit of dash or shortens the last segment.
\newcommand*\ulsbs@ccm@boundary[1]{\dimexpr\ulsbs@ccm@halfw
  +\ulsbs@ccm@arm*\numexpr#1\relax/\ulsbs@ccm@segments\relax}

% Build a single filled path with its elbow at the origin and legs along +x/+y.
% In dashed mode, the elbow includes the first dash of both legs; the remaining
% dashes are disconnected subpaths. No edges meet between separate paint calls.
\newcommand*\ulsbs@ccm@build{%
  \pgftransformreset
  \ulsbs@ccm@halfw=\dimexpr\ulChapterCornerMarkWidth/2\relax
  \ulsbs@ccm@reach=\ulChapterCornerMarkLength
  % A mark no longer than half its width consists of just the elbow.
  \ifdim\ulsbs@ccm@reach<\ulsbs@ccm@halfw
    \ulsbs@ccm@reach=\ulsbs@ccm@halfw
  \fi
  \ulsbs@ccm@segments=1
  \ifulChapterCornerMarkDashed
    \ifnum\ulChapterCornerMarkDashCount>0
      \ulsbs@ccm@segments=\numexpr2*\ulChapterCornerMarkDashCount+1\relax
    \fi
  \fi
  \ulsbs@ccm@arm=\dimexpr\ulsbs@ccm@reach-\ulsbs@ccm@halfw\relax
  \ulsbs@ccm@reach=\ulsbs@ccm@boundary{1}%
  \edef\ulsbs@ccm@h{\the\ulsbs@ccm@halfw}%
  \edef\ulsbs@ccm@r{\the\ulsbs@ccm@reach}%
  \pgfpathqmoveto{-\ulsbs@ccm@h}{-\ulsbs@ccm@h}%
  \pgfpathqlineto{ \ulsbs@ccm@r}{-\ulsbs@ccm@h}%
  \pgfpathqlineto{ \ulsbs@ccm@r}{ \ulsbs@ccm@h}%
  \pgfpathqlineto{ \ulsbs@ccm@h}{ \ulsbs@ccm@h}%
  \pgfpathqlineto{ \ulsbs@ccm@h}{ \ulsbs@ccm@r}%
  \pgfpathqlineto{-\ulsbs@ccm@h}{ \ulsbs@ccm@r}%
  \pgfpathclose
  \ifdim\ulsbs@ccm@arm>0pt
    \ulsbs@ccm@segment=2
    \@whilenum\ulsbs@ccm@segment<\ulsbs@ccm@segments\do{%
      \ulsbs@ccm@start=\ulsbs@ccm@boundary{\ulsbs@ccm@segment}%
      \ulsbs@ccm@end=\ulsbs@ccm@boundary{\ulsbs@ccm@segment+1}%
      \pgfpathrectanglecorners{\pgfqpoint{\ulsbs@ccm@start}{-\ulsbs@ccm@halfw}}%
        {\pgfqpoint{\ulsbs@ccm@end}{\ulsbs@ccm@halfw}}%
      \pgfpathrectanglecorners{\pgfqpoint{-\ulsbs@ccm@halfw}{\ulsbs@ccm@start}}%
        {\pgfqpoint{\ulsbs@ccm@halfw}{\ulsbs@ccm@end}}%
      \advance\ulsbs@ccm@segment by 2
    }%
  \fi
  \pgfusepathqfill
}

% Keep the most recent geometry of each style. Offset, orientation and chapter
% color are applied at use time, so neither page parity nor chapter changes
% require rebuilding the object.
\newcommand*\ulsbs@ccm@use[1]{%
  \edef\ulsbs@ccm@key{\the\ulChapterCornerMarkWidth/\the\ulChapterCornerMarkLength
    \ifulChapterCornerMarkDashed/\the\ulChapterCornerMarkDashCount\fi}%
  \expandafter\ifx\csname ulsbs@ccm@key@#1\endcsname\ulsbs@ccm@key
  \else
    \pgfdefobject{ulsbs@ccm@object@#1}%
      {\pgfqpoint{-.5\ulChapterCornerMarkWidth}{-.5\ulChapterCornerMarkWidth}}%
      {\pgfqpoint{\ulChapterCornerMarkLength}{\ulChapterCornerMarkLength}}%
      {\ulsbs@ccm@build}%
    \expandafter\xdef\csname ulsbs@ccm@key@#1\endcsname{\ulsbs@ccm@key}%
  \fi
  \pgfuseobject{ulsbs@ccm@object@#1}%
}

% Draws one chapter corner mark inside the caller's TikZ picture.
% #1 = anchor (north east / north west / south east / south west)
% #2 = x sign: + or -
% #3 = y sign: + or -
\newcommand*\ulChapterCornerMark[3]{%
  \pgfscope
    \pgftransformshift{\pgfpointadd{\pgfpointanchor{current page}{#1}}%
      {\pgfqpoint{#2\ulChapterCornerMarkOffset}{#3\ulChapterCornerMarkOffset}}}%
    \edef\ulsbs@ccm@xscale{\number\numexpr#21\relax}%
    \edef\ulsbs@ccm@yscale{\number\numexpr#31\relax}%
    \pgftransformcm{\ulsbs@ccm@xscale}{0}{0}{\ulsbs@ccm@yscale}{\pgfpointorigin}%
    \pgflowlevelsynccm
    \pgfsetfillcolor{\ulsbs@currentchaptercolor}%
    \ifulChapterCornerMarkDashed
      \ulsbs@ccm@use{dashed}%
    \else
      \ulsbs@ccm@use{solid}%
    \fi
  \endpgfscope
}

% Creates chapter corner marks on the outer corners of the page
% #1 = which side is outer?  R (odd pages) or L (even pages)
% Emits a zero-layout-space remembered TikZ overlay only when chapter colors are
% enabled. Caller supplies page parity; this macro does not advance page state.
\newcommand*\ulChapterCornerMarksOuter[1]{%
  \ifulUseChapterColors
    \begin{tikzpicture}[overlay,remember picture]
      \ifx#1R
        \ulChapterCornerMark{north east}{-}{-}%
        \ulChapterCornerMark{south east}{-}{+}%
      \else
        \ulChapterCornerMark{north west}{+}{-}%
        \ulChapterCornerMark{south west}{+}{+}%
      \fi
    \end{tikzpicture}%
  \fi
}


% Setup main header style
\fancypagestyle{ulsbs}[fancy]{%
  % NOTE: If fancydr warns that \headheight, \headsep or \footskip is too small,
  % adjust the value by defining \ulHeadHeight, \ulHeadSep or \ulFootSkip
  % commands before loading this package.
  \setlength{\footskip}{\ulFootSkip}
  \renewcommand{\headrulewidth}{0pt}
  \renewcommand{\footrulewidth}{0pt}
  \fancyhf{}
  %\fancyhead[LE,RO]{\ulPageNumberStyle\thepage} % page nr
  \fancyhead[CE]{\ulsbs@headerbooktitle}%    % even center: book title
  \fancyhead[CO]{\ulsbs@headerchapterandsectiontitles} % chapter title --- section title
  \fancyhead[RO]{%
    \ulPageNumberStyle\thepage%
    \ulChapterCornerMarksOuter{R}%
  }
  \fancyhead[LE]{%
    \ulPageNumberStyle\thepage%
    \ulChapterCornerMarksOuter{L}%
  }
}

% Setup style for TOC
\fancypagestyle{ulsbstoc}[ulsbs]{%
  \fancyhead[LO]{} % odd left: nothing
  \fancyhead[LE,RO]{%
    \ulPageNumberStyle\thepage%
  }
  \fancyhead[CO]{\rmfamily\small\itshape{\leftmark}} % odd center: chapter title: "CONTENTS"
}

\pagestyle{ulsbs} % set the default page style

% === Inline presentation boxes ===

% These boxes use tcolorbox package. https://ctan.org/pkg/tcolorbox
% They are used by giving color of the box as optional parameter, and the
% content as the mandatory one. If the color is omitted, current color is
% used by default. Note that the current color seems to be reset to black
% within the box, so to use it there, it needs to be saved and reset.
%
% Usage examples:
%   - \ulsbs@altchordidbox[red]{(Am)} % red box, black text
%   - \colorlet{ulsbs@curcol}{.}\ulsbs@altchordidbox{\color{ulsbs@curcol}(Am)} % current color for box and text

% Coloured box for language code.
\newtcbox{\ulsbs@langcodebox}[1][.]{%
  on line,arc=1pt,outer arc=.618pt,%
  colback=#1!6!white,colframe=#1!80!black,
  boxsep=0pt,%
  left=0.8pt,right=0.8pt,top=0.8pt,bottom=0.8pt,%
  boxrule=0.5pt%
}

% Coloured box for the alternate chords id (title).
\newtcbox{\ulsbs@altchordidbox}[1][.]{%
  on line,%
  valign=bottom,%
  arc=3pt,%outer arc=3pt,%
  colback=#1!5!white,colframe=#1!94!black,%
  boxsep=0pt,%
  left=1.4pt,right=1.4pt,top=0.1pt,bottom=0.1pt,%
  boxrule=0.3pt%
}

% Coloured box for the "Key: XX" text on the keyline of song prelude
\newtcbox{\ulsbs@mainkeybox}[1][.]{%
  on line,%
  valign=bottom,%
  arc=3pt,%outer arc=3pt,%
  colback=#1!8!white,colframe=#1!92!black,%
  boxsep=0pt,%
  left=1.4pt,right=1.4pt,top=0.1pt,bottom=0.1pt,%
  boxrule=0.5pt,%
  %equal height group=preludekeyheightgroup%
}

% Coloured box for the good keys, "Good: Am (Am--Bm)", text on the keyline
% of song prelude
\newtcbox{\ulsbs@goodkeysbox}[1][.]{%
  on line,%
  valign=bottom,%
  arc=3pt,%outer arc=3pt,%
  colback=#1!5!white,colframe=#1!94!black,%
  boxsep=0pt,%
  left=1.4pt,right=1.4pt,top=0.1pt,bottom=0.1pt,%
  boxrule=0.5pt,%
  %equal height group=preludekeyheightgroup%
}

% Coloured box for tags in the song prelude
\newtcbox{\ulsbs@tagsbox}[1][.]{%
  on line,%
  valign=bottom,%
  arc=1pt,%outer arc=2pt,%
  colback=#1!5!white,colframe=#1!94!black,%
  boxsep=0pt,%
  left=1.3pt,right=1.3pt,top=0.1pt,bottom=0.1pt,%
  boxrule=0.5pt,%
  %equal height group=preludekeyheightgroup%
}

% === Column break and stanza spacing policy ===

% Layout: column/page break penalties (<=-10000 forces a break).
% Check stanza/column break behavior when changing them. Pagination may evolve
% with styling; unchanged historical page counts are not a runtime requirement.
\ulSongPenalty=0     % between songs
\ulVerseVersePenalty=200  % between verse and verse
\ulBreakPenalty=50  % for \brk
\ulHardBreakPenalty=-500 % for \hardbrk
\interlinepenalty=1000 % between in-verse lines
\newcount\ulGluePenalty%
\ulGluePenalty=950 % between verses glued together with \glueverses (default here: 950)
\newcount\ulBlockPenalty%
\ulBlockPenalty=0 % between (and before the song's first) in-song blocks (default here: 0)
\newcount\ulTranslationVersePenalty%
\ulTranslationVersePenalty=900 % between translation verses (when \nextverse is called)

% Do not use verse numbers
\ulNoVerseNumbers
